\subsection{\texorpdfstring{同态 \(\operatorname{Ext}_{\mathbf{A}}(\mathbf{N}, \mathbf{P}) \otimes \operatorname{Ext}_{\mathbf{A}}(\mathbf{M}, \mathbf{N}) \to \operatorname{Ext}_{\mathbf{A}}(\mathbf{M}, \mathbf{P})\)}{同态 \textbackslash operatorname\{Ext\}\_\{\textbackslash mathbf\{A\}\}(\textbackslash mathbf\{N\}, \textbackslash mathbf\{P\}) \textbackslash otimes \textbackslash operatorname\{Ext\}\_\{\textbackslash mathbf\{A\}\}(\textbackslash mathbf\{M\}, \textbackslash mathbf\{N\}) \textbackslash 到 \textbackslash operatorname\{Ext\}\_\{\textbackslash mathbf\{A\}\}(\textbackslash mathbf\{M\}, \textbackslash mathbf\{P\})}}

设 M 和 N 是两个左 A-模。考虑拟同构 \(p_{\mathbf{M}}: \mathbf{L}(\mathbf{M}) \to \mathbf{M}\) , \(e_{\mathbf{M}}: \mathbf{M} \to \mathbf{l}(\mathbf{M})\) 与 \(e_{\mathbf{M}} \circ p_{\mathbf{M}}: \mathbf{L}(\mathbf{M}) \to \mathbf{l}(\mathbf{M})\) ；由 X，第 86 页，命题 4，我们由此推出一个双射同态

\[
a _ {\mathrm{M}, \mathrm{N}} = \mathrm{H} (\operatorname{Homgr} _ {\mathrm{A}} (e _ {\mathrm{M}} \circ p _ {\mathrm{M}}, 1)): \mathrm{H} (\operatorname{Homgr} _ {\mathrm{A}} (\mathrm{I} (\mathrm{M}), \mathrm{I} (\mathrm{N}))) \rightarrow \operatorname{Ext} _ {\mathrm{A}} (\mathrm{M}, \mathrm{N}).
\]

另外回忆（X，第 82 页）： \(\mathrm{H}^n (\mathrm{Homgr}_{\mathbf{A}}(\mathrm{I}(\mathbf{M}),\mathrm{I}(\mathbf{N})))\) 是从复形 I(M) 到复形 I(N) 的 \(n\) 次上升态射的同伦类集合。例如，若 \(f\in \mathrm{Hom}_{\mathbf{A}}(\mathbf{M},\mathbf{N})\) ，则 I(f) 的同伦类被 \(a_{\mathbf{M},\mathbf{N}}\) 映到 \(f\) .

若 P 是第三个左 A-模，根据 X，第 99 页，我们有一个典范的 k-同态

\[
\mathrm{H} (\operatorname{Homgr} _ {\mathbf {A}} (I (N), I (P))) \otimes_ {k} \mathrm{H} (\operatorname{Homgr} _ {\mathbf {A}} (I (M), I (N))) \rightarrow \mathrm{H} (\operatorname{Homgr} _ {\mathbf {A}} (I (M), I (P)))
\]

由此我们借助同构 \(a_{\mathbf{N},\mathbf{P}}, a_{\mathbf{M},\mathbf{N}}, a_{\mathbf{M},\mathbf{P}}\) 通过结构转移，推出一个 \(k\) -同态（复合同态）：

\[
c _ {\mathbf {M}, \mathbf {N}, \mathbf {P}}: \operatorname{Ext} _ {\mathbf {A}} (\mathbf {N}, \mathbf {P}) \otimes_ {k} \operatorname{Ext} _ {\mathbf {A}} (\mathbf {M}, \mathbf {N}) \rightarrow \operatorname{Ext} _ {\mathbf {A}} (\mathbf {M}, \mathbf {P}).
\]

这对应于一个 k-双线性映射。

\[
\operatorname{Ext} _ {\mathbf {A}} (\mathbf {N}, \mathbf {P}) \times \operatorname{Ext} _ {\mathbf {A}} (\mathbf {M}, \mathbf {N}) \rightarrow \operatorname{Ext} _ {\mathbf {A}} (\mathbf {M}, \mathbf {P})\tag{1}
\]

它分解为齐次分量

\[
\operatorname{Ext} _ {\mathbf {A}} ^ {i} (\mathbf {N}, \mathbf {P}) \times \operatorname{Ext} _ {\mathbf {A}} ^ {j} (\mathbf {M}, \mathbf {N}) \rightarrow \operatorname{Ext} _ {\mathbf {A}} ^ {i + j} (\mathbf {M}, \mathbf {P}).\tag{2}
\]

若 \(u \in \text{Ext}_A(\mathbf{N}, \mathbf{P})\) , \(v \in \text{Ext}_A(\mathbf{M}, \mathbf{N})\) ， \((u, v)\) 在 (1) 下的像称为 \(u\) 与 \(v\) 的复合积，记作 \(u \circ v\) 。如果 \(g\) （相应地， \(f\) ）是从 I(M) 到 I(N)（相应地，从 I(N) 到 I(P)）的 \(j\) （相应地， \(i\) ）次上升复形态射，其同伦类为 \(\overline{g}\) （相应地， \(\overline{f}\) ），则复合积 \(a_{\mathbf{N},\mathbf{P}}(\overline{f}) \circ a_{\mathbf{M},\mathbf{N}}(\overline{g})\) 是 \(a_{\mathbf{M},\mathbf{P}}\) 下从 I(M) 到 I(P) 的态射 \(f \circ g\) 的同伦类的像。

例 1. --- 如果 \(u \in \operatorname{Hom}_{\mathbf{A}}(\mathbf{N}, \mathbf{P})\) , \(v \in \operatorname{Hom}_{\mathbf{A}}(\mathbf{M}, \mathbf{N})\) , 则 \(u \circ v\) 是 \(u\) 与 \(v\) 的复合。

例2. --- 如果 \(u \in \mathrm{Hom}_{\mathbf{A}}(\mathbf{N}, \mathbf{P})\) , \(v \in \mathrm{Ext}_{\mathbf{A}}(\mathbf{M}, \mathbf{N})\) ，则

\[
u \circ v = \operatorname{Ext} _ {A} (1 _ {M}, u) (v) \in \operatorname{Ext} _ {A} (M, P);
\]

同样地，如果 \(u \in \operatorname{Ext}_{\mathbf{A}}(\mathbf{N}, \mathbf{P})\) , \(v \in \operatorname{Hom}_{\mathbf{A}}(\mathbf{M}, \mathbf{N})\) ，则

\[
u \circ v = \operatorname{Ext} _ {A} (v, 1 _ {p}) (u) \in \operatorname{Ext} _ {A} (M, P).
\]

这由 X，第 88 页的定义和注记得出。

若 Q、M、N、P 是四个左 A-模，且若

\[
u \in \operatorname{Ext} _ {\mathbf {A}} (\mathbf {N}, \mathrm{P}), \quad v \in \operatorname{Ext} _ {\mathbf {A}} (\mathbf {M}, \mathbf {N}), \quad w \in \operatorname{Ext} _ {\mathbf {A}} (\mathbf {Q}, \mathbf {M}),
\]

那么 \((u \circ v) \circ w = u \circ (v \circ w)\) ：复合积是结合的；因此，我们将不加括号地书写多个元素的复合。特别地，根据例 2：

例3. — 设 M、N、M'、N' 为四个左 A-模。若 \(u \in \operatorname{Ext}_{\mathbf{A}}(\mathbf{M}, \mathbf{N})\) , \(f \in \operatorname{Hom}_{\mathbf{A}}(\mathbf{M}', \mathbf{M})\) , \(g \in \operatorname{Hom}_{\mathbf{A}}(\mathbf{N}, \mathbf{N}')\) ，则

\[
\operatorname{Ext} _ {\mathbf {A}} (f, g) (u) = g \circ u \circ f \in \operatorname{Ext} _ {\mathbf {A}} (\mathbf {M} ^ {\prime}, \mathbf {N} ^ {\prime}).\tag{3}
\]

这给出了一个新的证明，即 \(k\) -双线性性，亦即映射 \((f, g) \to \operatorname{Ext}_{\mathbf{A}}(f, g)\) 的双线性性（X，第 88 页，命题 6）。

